\chapter{Windows, MinGW-w64, and Wine}
\label{ch:windows-wine}

Volume~I's Direct3D chapter described what the D3D9/D3D11/D3D12 backends prove; this chapter
describes the actual machinery that makes proving it possible without ever touching a
Windows machine --- a genuinely elaborate, carefully engineered dev loop that deserves its
own explanation.

\section{The cross-compilation toolchain}

\texttt{cmake/toolchains/mingw-w64.cmake} is a short, standard CMake cross-toolchain file:
it sets \texttt{CMAKE\_SYSTEM\_NAME} to \texttt{Windows}, locates the
\texttt{x86\_64-w64-mingw32-gcc}/\texttt{g++} toolchain (falling back to a 32-bit triple if
the 64-bit one is not found), and isolates the cross-root so CMake's own
\texttt{find\_package} calls never accidentally pick up a host Linux library. The invocation
is exactly what Chapter~\ref{ch:building-cna} in Volume~I already showed:
\texttt{-DCMAKE\_TOOLCHAIN\_FILE=cmake/toolchains/mingw-w64.cmake}, alongside whichever
backend option is being targeted.

\section{Three Wine prefixes, never one}

A single Wine prefix is not enough, and the project maintains three, each for a distinct
reason. One prefix carries a genuine Microsoft \texttt{d3dcompiler\_47.dll} (installed via
\texttt{winetricks}), needed because Wine's own \emph{builtin} \texttt{d3dcompiler\_47.dll} is
backed by an incomplete shader compiler --- it fails outright on a real alpha-test ternary
expression from the stock effects, and where it does compile something, the result can be
wildly wrong: a 72-bone skinning shader that should compile to about 2{,}160 bytes came out
of the builtin compiler at 1.4~MB, fully unrolled, because it cannot do relative addressing
at all. A second, 32-bit prefix carries a full XNA~4.0 installation --- real .NET Framework 4,
the real XNA~4.0 redistributable, and a real in-prefix \texttt{csc.exe} --- specifically so
the D3D9 oracle (\S\ref{sec:d3d9-oracle}) can compile and run genuine XNA~4.0 code, not a
substitute for it. A third prefix is the one CNA's own D3D11 test suite runs against day to
day, deliberately left alone rather than experimented on.

\section{The self-verifying gate: catching a silent fallback before it corrupts a result}

Every \texttt{scripts/run-wine-dxvk*.sh}/\texttt{run-wine-vkd3d.sh} runner shares one
critical property: none of them simply trust that DXVK or vkd3d-proton is actually in effect.
Each one captures the target program's own log output and greps it for a translation-layer-specific
marker --- a real \texttt{DXVK: <version>} line, or a real \texttt{vkd3d-proton -
applicationVersion: <version>} line --- and refuses to treat the run as valid otherwise,
exiting with a specific error naming the failure as ``a silent WineD3D fallback rather than a
real DXVK-backed Direct3D~11 run.'' This is not a hypothetical concern: without this gate, a
misconfigured prefix could silently render through Wine's own, much less accurate, built-in
Direct3D implementation instead of the real GPU-accelerated translation layer the whole
verification methodology depends on, and every pixel-test result from that run would be
comparing CNA against the wrong thing without anyone noticing.

\section{D3D12's extra wrinkle: Proton, not plain Wine}

Chapter~\ref{ch:direct3d-backends} in Volume~I already described the real, root-caused
architecture mismatch behind D3D12's swap-chain crash under plain Wine: Debian's own system
\texttt{dxgi.dll} cannot hand a command queue to vkd3d-proton's separately overridden
\texttt{d3d12.dll}, because vkd3d-proton's own \texttt{dxgi.dll} is only ABI-compatible with
the rest of its own matched build. The concrete fix is a dedicated script that launches
through Valve's own Proton launcher end to end, which supplies a self-consistent DLL set that
vkd3d-proton's override then layers onto correctly. This script carries a hard-coded safety
guard refusing to ever operate against a real personal Wine prefix --- added after a genuine
incident where an earlier version of the same script accidentally launched against a
developer's own \texttt{\textasciitilde/.wine} prefix and hung for nine and a half hours
under a first-run installation wizard.

\section{The D3D9 oracle, mechanically}
\label{sec:d3d9-oracle}

Chapter~\ref{ch:direct3d-backends}'s account of the D3D9 oracle can now be given its full
mechanical explanation. A tool compiles Microsoft's own, unmodified stock-effect \texttt{.fx}
source files to real Direct3D~9 bytecode via the genuine \texttt{d3dcompiler\_47.dll},
cross-built and run through Wine against the dedicated compiler prefix above; a separate
comparison step diffs the result against FNA's own shipped, Microsoft-compiled bytecode ---
61 of 66 compiled variants matched byte-for-byte, and the remaining five (all
\texttt{PixelLighting} vertex variants) were confirmed, not merely assumed, to differ only by
compiler version, after ruling out every optimization-flag combination as the actual cause.

The pixel-diff oracle itself works by rendering the identical declarative scene description
through two independent paths: a real, Wine-hosted XNA~4.0 program (compiled by the genuine
in-prefix \texttt{csc.exe}, run against the XNA-redistributable prefix, itself running
through DXVK rather than WineD3D --- a deliberate methodological choice, so that a pixel diff
never silently measures a driver difference instead of a real CNA difference) and CNA's own
cross-compiled D3D9 backend, run through the same DXVK-verified gate as everything else in
this chapter. A dedicated, mutation-tested comparison script does the actual pixel diff, at
a caller-specified tolerance --- confirmed, via a deliberately corrupted test image, to
correctly report exactly the number of differing pixels a known mutation introduces, not just
a pass/fail verdict.

\section{fna-reference: a different, lighter-weight comparison, and why it uses Mono instead}

Distinct from the rendering-capable D3D9 oracle above, \texttt{tools/fna-reference/} answers
a narrower question --- do CNA's enums, state-object presets, packed vector types, and
viewport math match FNA's own values exactly --- and does so without any Wine or Windows
involvement at all. It runs the real \texttt{FNA.dll} natively on Linux through Mono, dumping
a fixed set of non-rendering values (enum members, default state-object properties, packed
vector round-trips, viewport projection math) to JSON; a companion C\texttt{++} tool performs
the identical dump from CNA's own real implementation, entirely without constructing a
\cnaclass{GraphicsDevice} at all; and a small, dependency-free comparison script asserts every
FNA-side value has a matching CNA-side value, deliberately one-directional so that a
CNA-only, \texttt{NOXNA}-tagged extension is never flagged as a mismatch. This lighter-weight
tool exists specifically because it needs no live rendering device --- extending it to cover
anything that does (a stock effect's actual rendered output, for instance) was evaluated and
explicitly deferred as a substantially larger undertaking, left for the D3D9 oracle to cover
instead. One real, if narrow, divergence this tool did catch: \cnaclass{IndexElementSize}'s
\texttt{SixteenBits}/\texttt{ThirtyTwoBits} values did not match FNA's own \texttt{0}/\texttt{1}
numbering at all --- confirmed and fixed once this comparison surfaced it.
